\documentclass[11pt]{article}

\usepackage[margin=1in]{geometry}
\usepackage{amsmath,amssymb,amsthm,mathtools}
\usepackage[hidelinks]{hyperref}

\newtheorem{theorem}{Theorem}[section]
\newtheorem{lemma}[theorem]{Lemma}
\newtheorem{proposition}[theorem]{Proposition}
\newtheorem{corollary}[theorem]{Corollary}
\theoremstyle{remark}
\newtheorem{remark}[theorem]{Remark}

\newcommand{\one}{\mathbf 1}
\newcommand{\tr}{\operatorname{tr}}
\newcommand{\DNN}{\mathrm{DNN}}

\title{The Optimized Liu--Tang--Zhang Constant for Cactus Graphs}
\author{}
\date{25 July 2026}

\begin{document}
\maketitle

\begin{abstract}
For a connected graph $G$, we optimize the doubly nonnegative functional
underlying the square-energy theorem of Liu, Tang, and Zhang.  If $G$ is a
cactus with $b$ bridge blocks and cyclic blocks of lengths
$\ell_1,\ldots,\ell_r$, then its exact optimized constant is
\[
 \kappa(G)=b+\sum_{j=1}^r\kappa(C_{\ell_j}),
 \qquad
 \kappa(C_\ell)=
 \begin{cases}
  \ell,&\ell\text{ even},\\
  \ell\sec^2(\pi/(2\ell)),&\ell\text{ odd}.
 \end{cases}
\]
 The proof combines a direct Fourier optimization of the
 Liu--Tang--Zhang functional on a cycle with exact additivity under a
 one-vertex sum.
The latter requires both directions: folding and a Gram split give the upper
bound, while orthogonal Gram gluing with the correct relative scaling gives
the lower bound.  Applied to the negative spectral part of the adjacency
matrix, the result yields $s^-(G)\leq\kappa(G)$.  For bicyclic cacti this
settles the positive square-energy conjecture of
Akbari--Kumar--Mohar--Pragada--Zhang except possibly when the unordered pair
of cyclic block lengths is $\{5,5\}$ or $\{3,4k+1\}$, $k\geq1$.  These are
unresolved families here, not asserted exceptions or counterexamples.
\end{abstract}

\section{Introduction}

All graphs are finite, simple, and undirected.  If the adjacency eigenvalues
of $G$ are $\lambda_1,\ldots,\lambda_n$, write
\[
 s^+(G)=\sum_{\lambda_i>0}\lambda_i^2,
 \qquad
 s^-(G)=\sum_{\lambda_i<0}\lambda_i^2.
\]
These quantities were introduced in connection with spectral bounds for the
chromatic number by Elphick, Farber, Goldberg, and Wocjan~\cite{EFGW}.  Their
conjecture $\min\{s^+,s^-\}\geq n-1$ for connected graphs was recently proved
by Liu, Tang, and Zhang~\cite{LTZ}.  The central device in their proof is an
inequality for doubly nonnegative matrices, meaning real symmetric matrices
that are both positive semidefinite and entrywise nonnegative.

We optimize that device graph by graph.  For a connected graph with at least
one edge, set
\begin{equation}\label{eq:kappa-definition}
 \kappa(G)=
 \sup_{\substack{M\succeq0,\ M\geq0\\T(M)>0}}
 \frac{4S_G(M)^2}{T(M)},
 \qquad
 S_G(M)=\sum_{uv\in E(G)}\sqrt{M_{uv}},
 \quad T(M)=\one^TM\one.
\end{equation}
Thus $M$ ranges over the doubly nonnegative cone.  The universal
Liu--Tang--Zhang inequality says
\[
 \kappa(G)\leq 2|E(G)|-|V(G)|+1.
\]
Our purpose is not to replace their framework but to compute its sharp
constant on cacti.

A \emph{cactus} is a connected graph in which every block is a bridge or a
cycle.  Equivalently, any two cycles have at most one vertex in common.  Our
main theorem is the following.

\begin{theorem}\label{thm:main}
Let $G$ be a connected cactus with $b$ bridge blocks and cyclic blocks
$C_{\ell_1},\ldots,C_{\ell_r}$.  Then
\begin{equation}\label{eq:main}
 \kappa(G)=b+\sum_{j=1}^r\kappa(C_{\ell_j}),
\end{equation}
where
\begin{equation}\label{eq:cycle-value}
 \kappa(C_\ell)=
 \begin{cases}
  \ell,&\ell\text{ is even},\\[2mm]
  \ell\sec^2\!\left(\dfrac{\pi}{2\ell}\right),
       &\ell\text{ is odd}.
 \end{cases}
\end{equation}
\end{theorem}

The word ``sharp'' in Theorem~\ref{thm:main} refers to the optimized DNN
certificate~\eqref{eq:kappa-definition}.  It does not say that the spectral
bound $s^-(G)\leq\kappa(G)$ is always an equality.

Akbari, Kumar, Mohar, Pragada, and Zhang conjectured that every connected
$n$-vertex graph with at least $n+1$ edges satisfies $s^+(G)\geq n$
\cite[Conjecture~1.2]{AKMPZ}.  The exact cactus constant proves this for a
large class of bicyclic cacti.  The remaining common-phase cases are supplied
by a direct packing-two Sachs argument: $s^+>s^-$ when every cycle has length
$3\pmod4$ and at most two cycles can be pairwise vertex-disjoint.  Combining
the two methods gives an exact account of what is proved in this class; see
Theorem~\ref{thm:classification}.

\section{The exact cyclic certificate}

We evaluate the cyclic instance of the Liu--Tang--Zhang functional directly.
Given a DNN matrix $M$, fold each
nonedge entry onto the corresponding two diagonal entries:
\begin{equation}\label{eq:folding}
 \mathcal F_G(M)=M+
 \sum_{\substack{\{u,v\}\notin E(G)\\u\ne v}}
 M_{uv}(e_u-e_v)(e_u-e_v)^T.
\end{equation}
Each summand is positive semidefinite and entrywise nonnegative off the two
cancelled positions.  Consequently $\mathcal F_G(M)$ is DNN, vanishes on
nonedges, has the same edge entries as $M$, and has the same value of $T$.
This is the folding step in the LTZ argument.

\begin{proposition}\label{prop:cycle}
For every $\ell\geq3$,
\[
 \kappa(C_\ell)=
 \begin{cases}
  \ell,&\ell\text{ even},\\[1mm]
  \ell\sec^2(\pi/(2\ell)),&\ell\text{ odd}.
 \end{cases}
\]
\end{proposition}

\begin{proof}
 Fold $M$ by~\eqref{eq:folding}, and then average it over all cyclic
 rotations.  The resulting matrix $\overline M$ is circulant and DNN, while
 $T(\overline M)=T(M)$.  If $t$ is its common cycle-edge entry, concavity of
 the square root gives
 \[
  S_{C_\ell}(\overline M)
  =\ell\sqrt{\frac1\ell\sum_{e\in E(C_\ell)}M_e}
  \geq S_{C_\ell}(M).
 \]
 Write $r$ for the common row sum of $\overline M$.  Then
 \[
  T(\overline M)=\ell r,
  \qquad
  \frac{4S(\overline M)^2}{T(\overline M)}=\frac{4\ell t}{r}.
 \]

 Suppose first that $\ell=2m$.  Let $x_d\geq0$ be the circulant entry at
 cyclic distance $d$, so $x_1=t$.  The Fourier eigenvalue belonging to the
 alternating vector is
 \[
  \lambda_m=x_0+2\sum_{d=1}^{m-1}(-1)^d x_d+(-1)^m x_m\geq0.
 \]
 Comparing it with the row sum gives
 \[
  r-\lambda_m
  =4\sum_{\substack{1\leq d<m\\d\text{ odd}}}x_d
   +2\mathbf1_{\{m\text{ odd}\}}x_m
  \geq4t.
 \]
 Hence $r\geq4t$ and the objective is at most $\ell$.

 Now let $\ell=2m+1$ and put $c=\cos(\pi/\ell)$.  The Fourier eigenvalue at
 frequency $m$ is nonnegative.  Therefore
 \begin{align*}
  r-\lambda_m
  &=2\sum_{d=1}^{m}
    \left(1-\cos\frac{2\pi md}{\ell}\right)x_d\\
  &\geq2(1+c)x_1=2(1+c)t,
 \end{align*}
 where only the $d=1$ term was retained and all other terms are
 nonnegative.  Thus $r\geq2(1+c)t$, and the objective is at most
 \[
  \frac{2\ell}{1+c}=\ell\sec^2\!\left(\frac{\pi}{2\ell}\right).
 \]

For the reverse inequality let $A=A(C_\ell)$, choose $x>0$, and put
\begin{equation}\label{eq:cycle-primal}
 M=x(2cI+A),
 \qquad
 c=
 \begin{cases}
  1,&\ell\text{ even},\\
  \cos(\pi/\ell),&\ell\text{ odd}.
 \end{cases}
\end{equation}
The least eigenvalue of $A(C_\ell)$ is $-2$ for even $\ell$ and
$-2\cos(\pi/\ell)$ for odd $\ell$, so $M$ is positive semidefinite.  It is
also entrywise nonnegative.  Every edge entry is $x$, and
\[
 S=\ell\sqrt{x},
 \qquad T=2x\ell(c+1).
\]
Thus the objective is $2\ell/(c+1)$, which equals $\ell$ in the even case and
$\ell\sec^2(\pi/(2\ell))$ in the odd case.  This matches the certificate.
\end{proof}

\section{Exact one-vertex-sum additivity}

Write $G=G_1\vee_zG_2$ when $G_1$ and $G_2$ are connected edge-containing
graphs whose intersection is exactly the vertex $z$, with no edges between
$V(G_1)\setminus\{z\}$ and $V(G_2)\setminus\{z\}$.

\begin{theorem}[One-vertex-sum additivity]\label{thm:additivity}
For $G=G_1\vee_zG_2$,
\[
 \kappa(G)=\kappa(G_1)+\kappa(G_2).
\]
\end{theorem}

\begin{proof}
Put $\kappa_i=\kappa(G_i)$.  For the upper bound, first fold a feasible $M$
on $G$.  Let $g_v$ be Gram vectors for the resulting edge-supported matrix,
and set
\[
 \mathcal U=\operatorname{span}\{g_u:u\in V(G_1)\setminus\{z\}\},
 \qquad
 \mathcal W=\operatorname{span}\{g_w:w\in V(G_2)\setminus\{z\}\}.
\]
The nonedge zeroes imply $\mathcal U\perp\mathcal W$.  Decompose
$g_z=p_U+p_W+p_0$, where $p_U\in\mathcal U$, $p_W\in\mathcal W$, and
$p_0\perp\mathcal U+\mathcal W$.  On $G_1$ use the original vectors away
from $z$ and the cut vector $p_U+p_0$ at $z$; on $G_2$ use the original
vectors away from $z$ and $p_W$ at $z$.  Their Gram matrices $M_1,M_2$ are
DNN, all edge inner products are preserved, and orthogonality gives
\begin{equation}\label{eq:split-identities}
 S_G(M)=S_{G_1}(M_1)+S_{G_2}(M_2),
 \qquad T(M)=T(M_1)+T(M_2).
\end{equation}
Writing $S_i=S_{G_i}(M_i)$ and $T_i=T(M_i)$, we obtain
\[
 2(S_1+S_2)
 \leq\sqrt{\kappa_1T_1}+\sqrt{\kappa_2T_2}
 \leq\sqrt{(\kappa_1+\kappa_2)(T_1+T_2)}.
\]
This proves $\kappa(G)\leq\kappa_1+\kappa_2$.

For the lower bound, choose folded, edge-supported near-optimizers on the two
graphs and represent them in orthogonal Euclidean spaces.  Glue them by
replacing the two vectors at $z$ by their orthogonal direct sum; keep every
other vector in its own summand.  The resulting Gram matrix is DNN on $G$,
has no cross-side entries, and again satisfies~\eqref{eq:split-identities}.
Scale the two input matrices so that
\begin{equation}\label{eq:correct-scaling}
 \frac{T_1}{\kappa_1}=\frac{T_2}{\kappa_2}.
\end{equation}
For exact optimizers this makes both Cauchy--Schwarz inequalities above
equalities and gives objective value $\kappa_1+\kappa_2$.  Applying the same
scaling to arbitrarily close near-optimizers and taking a limit proves the
claim without assuming that either supremum is attained.
\end{proof}

\begin{proposition}\label{prop:bridge}
$\kappa(K_2)=1$.
\end{proposition}

\begin{proof}
For a DNN matrix $\left(\begin{smallmatrix}a&c\\c&d\end{smallmatrix}\right)$,
positive semidefiniteness gives $a+d\geq2c$, and hence
$T=a+d+2c\geq4c$.  Therefore the objective is at most one.  Equality holds
for any positive multiple of the all-ones matrix $J_2$.
\end{proof}

Every cactus is obtained by iterated one-vertex sums of its blocks.  Applying
Theorem~\ref{thm:additivity}, Proposition~\ref{prop:bridge}, and
Proposition~\ref{prop:cycle} proves Theorem~\ref{thm:main}.

\section{The spectral consequence}

Let $A=A(G)$ and write its positive and negative spectral parts as
\[
 A=P-B,
 \qquad P=A_+\succeq0,
 \quad B=A_-\succeq0,
 \quad PB=0.
\]
Thus $s^-(G)=\tr B^2$.  The next consequence is valid for every connected
graph, with its own optimized constant.

\begin{theorem}\label{thm:spectral}
If $G$ is connected and has at least one edge, then
\[
 s^-(G)\leq\kappa(G).
\]
\end{theorem}

\begin{proof}
The Schur product theorem makes $M=B\circ B$ doubly nonnegative.  Since
$A$ has two unit entries for every edge and $PB=0$,
\[
 \tr(AB)=\tr((P-B)B)=-\tr B^2=-s^-(G),
\]
and therefore
\[
 \frac{s^-(G)}2=-\sum_{uv\in E(G)}B_{uv}
 \leq\sum_{uv\in E(G)}|B_{uv}|.
\]
Moreover,
$T(B\circ B)=\sum_{u,v}B_{uv}^2=\tr B^2=s^-(G)$.  The defining DNN
inequality for $\kappa(G)$ gives the chain
\begin{equation}\label{eq:spectral-chain}
 s^-(G)^2
 \leq4\left(\sum_{uv\in E(G)}|B_{uv}|\right)^2
 \leq\kappa(G)\,\one^T(B\circ B)\one
 =\kappa(G)s^-(G).
\end{equation}
Since $G$ has an edge, $s^-(G)>0$, and division proves the result.
\end{proof}

\begin{remark}
Even when the cactus formula for $\kappa(G)$ is attained by a DNN matrix,
equality need not hold in the first inequality of~\eqref{eq:spectral-chain},
nor must $B\circ B$ optimize the DNN functional.  This is why DNN sharpness
must not be interpreted as spectral equality.
\end{remark}

\section{Bicyclic cacti}

Define the odd-cycle excess
\begin{equation}\label{eq:epsilon}
 \epsilon_\ell=
 \begin{cases}
  0,&\ell\text{ even},\\
  \ell\tan^2(\pi/(2\ell)),&\ell\text{ odd}.
 \end{cases}
\end{equation}
Then $\kappa(C_\ell)=\ell+\epsilon_\ell$.

\begin{lemma}[Exact trigonometric comparison]\label{lem:trig}
The sequence $\epsilon_\ell$ is strictly decreasing through odd integers
$\ell\geq3$, and
\[
 \epsilon_3=1,
 \qquad
 \epsilon_5=5-2\sqrt5,
 \qquad
 \epsilon_5+\epsilon_7<1.
\]
\end{lemma}

\begin{proof}
Put $t=\pi/(2x)$.  Apart from the positive constant $\pi/2$,
$x\tan^2(\pi/(2x))$ is $\tan^2t/t$.  Its derivative with respect to $t$ is
positive because
\[
 2t\tan t\sec^2t-\tan^2t>0
 \quad\Longleftrightarrow\quad
 2t>\sin t\cos t,
\]
 which holds on the relevant interval $0<t\leq\pi/6$.  Since $t$ decreases with $x$, the required
monotonicity follows.  The first identity is immediate, and
$\tan^2(\pi/10)=1-2\sqrt5/5$ gives the second.

For a fully strict comparison in the last assertion, $\pi<22/7$ and
$\cos x>1-x^2/2$ give
\[
 \cos(\pi/7)>1-\frac12\left(\frac{22}{49}\right)^2>\frac78.
\]
Hence
\[
 \epsilon_7
 =7\frac{1-\cos(\pi/7)}{1+\cos(\pi/7)}
 <\frac7{15}<2\sqrt5-4=1-\epsilon_5.
\]
For the middle strict inequality, $7/15<2\sqrt5-4$ is equivalent to
$67/30<\sqrt5$, whose square follows from $4489<4500$.
\end{proof}

Let $G$ now be a connected bicyclic cactus, and let $p,q$ be the lengths of
its two cyclic blocks.  Since $|E(G)|=|V(G)|+1$, if $n=|V(G)|$ and $b$ is
the number of bridges, then $b+p+q=n+1$.  Theorem~\ref{thm:main} and
Theorem~\ref{thm:spectral} therefore give
\begin{equation}\label{eq:bicyclic-negative}
 s^-(G)\leq n+1+\epsilon_p+\epsilon_q.
\end{equation}
Since $s^+(G)+s^-(G)=\tr A^2=2|E(G)|=2n+2$, we obtain the following direct
criterion.

\begin{corollary}\label{cor:dnn-bicyclic}
If $\epsilon_p+\epsilon_q\leq1$, then $s^+(G)\geq n$.  In particular, this
holds whenever at least one of $p,q$ is even, and also whenever $p,q$ are odd
with $\min\{p,q\}\geq5$ and $\{p,q\}\ne\{5,5\}$.
\end{corollary}

\begin{proof}
Equation~\eqref{eq:bicyclic-negative} gives
$s^+(G)\geq n+1-\epsilon_p-\epsilon_q\geq n$.  Even lengths have zero
excess, while every odd excess is at most $\epsilon_3=1$, proving all cases
with an even length.  For two odd lengths at least five, the only pair not
bounded by $\epsilon_5+\epsilon_7<1$ and monotonicity is $\{5,5\}$.
\end{proof}

The DNN estimate does not settle pairs containing a triangle and another odd
cycle, because $\epsilon_3=1$ and the second excess is positive.  When the
second cycle also has length $3\pmod4$, the independent phase method does
more than is needed.

\begin{proposition}[Packing-two consequence]
\label{prop:packing}
If $p\equiv q\equiv3\pmod4$, then
\[
 s^+(G)>s^-(G),
 \qquad\text{and consequently}\qquad
 s^+(G)>|E(G)|=n+1.
\]
\end{proposition}

\begin{proof}
 Let
 \[
  Z_H(t)=\sum_jm_j(H)t^{|V(H)|-2j}>0\qquad(t>0)
 \]
 be the signless matching polynomial.  Grouping the Sachs expansion by its
 collection $\mathcal C$ of vertex-disjoint cycle components gives
 \[
  i^{-|V(G)|}\phi_G(it)
  =\sum_{\mathcal C}
   \prod_{C\in\mathcal C}(-2i^{-|C|})
   Z_{G-V(\mathcal C)}(t).
 \]
 A bicyclic cactus has only two cyclic blocks, so $|\mathcal C|\leq2$.  If
 both lengths are $3\pmod4$, every singleton cycle has phase $-2i$, while
 the empty and two-cycle terms are real.  Hence the normalized characteristic
 polynomial lies strictly in the open lower half-plane for every $t>0$.
 Its continuous argument
 \[
  \Theta_G(t)=\sum_j\arctan(\lambda_j/t)
 \]
 is therefore negative.  The signed Coulson identity
 \[
  s^+(G)-s^-(G)=-\frac4\pi\int_0^\infty t\Theta_G(t)\,dt
 \]
 yields $s^+>s^-$.  The energy identity then turns this into
 $s^+>|E(G)|$.
\end{proof}

We can now state the broad classification supplied by the two methods.

\begin{theorem}[Bicyclic cactus classification]\label{thm:classification}
Let $G$ be a connected bicyclic cactus on $n$ vertices, with unordered pair
of cyclic block lengths $\{p,q\}$.  Then $s^+(G)\geq n$ in all cases except
possibly
\begin{equation}\label{eq:unresolved}
 \{p,q\}=\{5,5\}
 \qquad\text{or}\qquad
 \{p,q\}=\{3,4k+1\}\quad(k\geq1).
\end{equation}
More explicitly:
\begin{enumerate}
\item if at least one length is even, the optimized DNN bound proves the
      conjecture;
\item if both lengths are odd and at least five, the optimized DNN bound
      proves it unless both are five;
\item if both lengths are $3\pmod4$, the packing-two Sachs argument proves the
      stronger inequality $s^+(G)>|E(G)|=n+1$;
\item after these cases, exactly the pairs in~\eqref{eq:unresolved} remain.
\end{enumerate}
In particular, the conjecture of
Akbari--Kumar--Mohar--Pragada--Zhang holds for every bicyclic cactus outside
the two displayed families.
\end{theorem}

\begin{proof}
Corollary~\ref{cor:dnn-bicyclic} handles all even cases and all odd pairs with
both lengths at least five except $\{5,5\}$.  It remains to inspect pairs
containing a triangle.  If the other odd length is $3\pmod4$, including
another triangle, Proposition~\ref{prop:packing} applies.  The only untreated
odd partner is therefore $4k+1$ with $k\geq1$.  This gives precisely
\eqref{eq:unresolved}.
\end{proof}

\begin{remark}[Status of the residual families]
The word ``possibly'' is essential.  Neither $\{5,5\}$ nor
$\{3,4k+1\}$ is claimed to be an exception to the AKMPZ conjecture, and no
counterexample is asserted.  They are exactly the families not resolved by
the optimized cactus DNN certificate together with the cited packing-two
theorem.  Arbitrary bridge paths, shared cut vertices, and attached trees are
already included in this statement.
\end{remark}

\begin{thebibliography}{99}

\bibitem{AKMPZ}
S. Akbari, H. Kumar, B. Mohar, S. Pragada, and S. Zhang,
\emph{Refinement of a conjecture on positive square energy of graphs},
arXiv:2506.07264, 2025.

\bibitem{EFGW}
C. Elphick, M. Farber, F. Goldberg, and P. Wocjan,
\emph{Conjectured bounds for the sum of squares of positive eigenvalues of a
graph},
Discrete Math. 339 (2016), 2215--2223.

\bibitem{LTZ}
Y. Liu, Q. Tang, and S. Zhang,
\emph{The positive and negative square-energy conjecture},
arXiv:2607.18031, 2026.

\end{thebibliography}

\end{document}
